\section{Decision Problem and Frozen Protocol}
\label{sec:protocol}

\subsection{Graph-vs-MLP Decision}

The unit of evaluation is a dataset--seed pair $u$. Within each unit, validation performance selects one feature-only MLP and one graph model from GCN, GAT, GraphSAGE, H2GCN, LINKX, and GPR-GNN. Every architecture receives the same four-trial tuning budget, and exact validation ties are broken by model identifier. Let $a_u^{\mathrm{M}}$ and $a_u^{\mathrm{G}}$ denote the held-out test accuracies of the selected MLP and graph model. The evaluator reads each selected test result once; no diagnostic can affect training or hyperparameter selection.

We ask whether the graph family provides a practically material gain rather than whether it wins by an arbitrarily small amount. With the frozen margin $\delta=0.01$, the target action is
\begin{equation}
y_u =
\begin{cases}
\mathrm{G}, & a_u^{\mathrm{G}}-a_u^{\mathrm{M}}>\delta,\\
\mathrm{M}, & a_u^{\mathrm{G}}-a_u^{\mathrm{M}}\leq\delta.
\end{cases}
\label{eq:decision_target}
\end{equation}
The same tie policy applies to every evaluated rule. This formulation makes the decision asymmetric by design: a graph model must exceed the MLP by more than one percentage point to justify the graph action.

\subsection{Matched Training and Selection}

Table~\ref{tab:frozen_training} records the common training budget from the frozen machine-readable specification. The implementation uses PyTorch~\cite{paszke2019pytorch}, and each trial minimizes cross-entropy with Adam. Early stopping retains the checkpoint with the highest validation accuracy, then the lowest validation loss and earliest epoch; training stops after 100 stale epochs or 500 epochs in total. The four trial configurations are selected independently within every architecture, and the chosen checkpoint then receives one held-out test evaluation. Thus no model family obtains a larger search budget or repeated access to the test split.

\begin{table}[t]
\centering
\caption{Frozen Matched Training Configuration}
\label{tab:frozen_training}
\small
\begin{tabular}{@{}lp{0.64\linewidth}@{}}
\toprule
Component & Frozen setting \\
\midrule
Representation width & 64 hidden channels for every architecture. \\
Optimization & Adam, cross-entropy loss, weight decay $5\times 10^{-4}$. \\
Training horizon & At most 500 epochs; patience 100 under the validation checkpoint rule. \\
Four-trial grid & Learning rate in $\{0.01,0.005\}$ crossed with dropout in $\{0.5,0.7\}$. \\
Model-specific constants & GAT: 4 heads; H2GCN: 2 propagation layers with strict two-hop edges; LINKX: 2 outer, 1 edge, and 2 node layers; GPR-GNN: 10 propagation steps, $\alpha=0.1$, PPR initialization. \\
Selection and test & Validation accuracy descending, validation loss ascending, trial identifier ascending; one held-out test evaluation. \\
\bottomrule
\end{tabular}
\end{table}

The split is stratified 60/20/20 for train/validation/test and is regenerated deterministically for seeds 0--9. Every unit stores the source commit, canonical configuration digest, processed-data checksums, split identifier, four trial records, selected checkpoint, and environment provenance. The v1 and v2 JSON specifications differ only in the pre-outcome eligibility amendment described in Section~\ref{sec:prospective}; the later Squirrel amendment changes only the execution timeout.

AI-assisted tools supported software review, test development, reproducibility checks, and editorial work but did not generate experimental data or make reporting decisions. All reported experiments were executed under the frozen protocol, and the author independently verified the code, artifacts, analyses, and manuscript claims.

\subsection{Decision-Time Information and Policies}

Full node features and graph structure are available because they are model inputs. Label-dependent graph statistics may use training labels only. In particular, edge homophily uses edges whose endpoints are both labeled for training, and the two-hop statistic uses training-labeled endpoint pairs. Validation accuracy is available because model tuning precedes the decision. Test labels, test accuracy, and any statistic computed from test labels are forbidden diagnostic inputs; the evaluator rejects records that violate this boundary.

Table~\ref{tab:diagnostic_policies} lists the nine frozen policies. The thresholds are historical or deliberately simple sanity baselines; none is fitted to the prospective outcomes. The combined policy chooses graph when training-only edge homophily is at least $0.55$. Below that threshold, it chooses MLP when MLP validation accuracy is at least $0.40$ and otherwise abstains. Missing diagnostic inputs also cause abstention. For full-set evaluation, every abstention uses the same predeclared MLP fallback, while coverage remains separately visible.

\begin{table}[t]
\centering
\caption{Frozen Graph-vs-MLP Decision Policies}
\label{tab:diagnostic_policies}
\small
\begin{tabular}{@{}lp{0.67\linewidth}@{}}
\toprule
Policy & Decision rule \\
\midrule
Always-MLP & Select MLP for every unit. \\
Always-graph & Select the validation-chosen graph model for every unit. \\
Random 50/50 & Use the analytic expectation of equally likely graph and MLP actions. \\
Homophily-only & Select graph when training-only edge homophily is at least $0.50$. \\
Degree-only & Select graph when mean degree is at least $10$. \\
Homophily+degree & Threshold the equal-weight sum of centered homophily and log-degree scores at zero. \\
Two-hop-only & Select graph when training-only $h_2-h_1>0.05$. \\
Combined & Apply the historical homophily/MLP-validation rule described in the text; otherwise abstain. \\
Validation selection & Select graph when its validation advantage exceeds $\delta$. \\
\bottomrule
\end{tabular}
\end{table}

\subsection{Outcomes and Inference}

Selection accuracy measures agreement with Eq.~\eqref{eq:decision_target} after the common fallback, and coverage is the fraction of non-abstaining decisions. Because an incorrect action can have negligible or substantial cost, the primary loss is oracle-family regret,
\begin{equation}
r_u(\widehat{y})=\max\{a_u^{\mathrm{G}},a_u^{\mathrm{M}}\}-a_u^{\widehat{y}},
\label{eq:oracle_regret}
\end{equation}
where $a_u^{\widehat{y}}$ is the test accuracy of the chosen family after fallback. We report mean full-set regret for all units and covered-set regret for non-abstaining units. Risk--coverage curves order decisions by each policy's frozen confidence score.

Seeds remain paired within a dataset. We first average regret over seeds for each dataset, then use the dataset means for percentile-bootstrap intervals and two-sided sign-flip tests. Holm correction controls the family-wise error rate across the eight predeclared comparisons with the combined policy. Non-significance is not interpreted as equivalence.

The prospective specification, immutable per-unit schema, analysis code, bootstrap seed, permutation seed, and paper-level decision rule were fixed before assembling or evaluating the new comparative outcomes. Section~\ref{sec:prospective} reports two outcome-independent protocol amendments and keeps the earlier feasibility artifacts outside the final analysis.
